\documentclass{beamer}

\usepackage{beamerthemeSingapore}
\usepackage[brazil]{babel}
\usepackage[T1]{fontenc}
\usepackage[utf8]{inputenc}
\usepackage{graphicx}
\usepackage{url}
\usepackage{xcolor}
\usepackage{booktabs}
\usepackage{ragged2e}
\usepackage{etoolbox}
\usepackage{minted}
\usepackage{tikz}
\usetikzlibrary{positioning, arrows.meta, shapes.geometric}

\newcommand{\tcb}[1]{\textcolor{blue}{#1}}
\newcommand{\tcr}[1]{\textcolor{red}{#1}}
\newcommand{\tcg}[1]{\textcolor{green!50!black}{#1}}
\newcommand{\ddash}{-\mbox{}-}

\setminted{
  fontsize=\footnotesize,
  breaklines=true,
  breakanywhere=true,
  frame=lines,
  framesep=2mm,
  baselinestretch=1.05,
  tabsize=2
}

\apptocmd{\frame}{}{\justifying}{}

\setbeamertemplate{navigation symbols}{}

\title{05 - CSS3: fundamentos}
\author{Prof. Alex Kutzke}
\date{Agosto de 2026}

\begin{document}

% ==========================================
\begin{frame}
	\begin{center}
		\large{Setor de Educação Profissional e Tecnológica}\\
		Tecnologia em Análise e Desenvolvimento de Sistemas\\
		\vspace{1cm}
		\small{DS122 - Desenvolvimento de Aplicações Web 1}
	\end{center}
	\titlepage
\end{frame}

% ==========================================
\begin{frame}{Objetivos da aula}
	Ao final desta aula você deve ser capaz de:
	\begin{enumerate}
		\item Ligar uma folha de estilo externa ao documento;
		\item Escrever seletores de tipo, classe, \texttt{id}, descendente, atributo
		      e pseudo-classe;
		\item Prever qual declaração vence um conflito;
		\item Calcular a largura que um elemento ocupa na tela;
		\item Escolher entre \texttt{px}, \texttt{rem}, \texttt{em} e \texttt{\%};
		\item Declarar cores e verificar o contraste com o fundo;
		\item Definir a tipografia da página;
		\item Descobrir no inspetor por que uma regra não foi aplicada.
	\end{enumerate}
\end{frame}

% ==========================================
\begin{frame}{Roteiro}
	\tableofcontents
\end{frame}

% ==========================================
\section{Fundamentos}
% ==========================================

\begin{frame}{A página que já tinha aparência}
	A página da aula passada não tem CSS nenhum. Mesmo assim:

	\begin{itemize}
		\item o \texttt{<h1>} aparece grande e em negrito;
		\item os itens de \texttt{<ul>} aparecem com marcador;
		\item o link aparece azul e sublinhado.
	\end{itemize}

	\vspace{0.4cm}
	Isso vem da \textbf{folha de estilo do agente de usuário}, que o navegador
	aplica a todo documento para que ele seja legível sem CSS.

	\vspace{0.3cm}
	No inspetor (\texttt{F12}), as regras dela aparecem marcadas como
	\emph{user agent stylesheet}.
\end{frame}

\begin{frame}{Escrever CSS é substituir aquelas decisões}
	\begin{block}{O que muda com o CSS}
		O navegador continua aplicando a folha dele para tudo o que você não disser.
		Cada regra sua substitui uma decisão daquela folha.
	\end{block}

	\vspace{0.5cm}
	CSS: \emph{Cascading Style Sheets}. A linguagem descreve a
	\textbf{apresentação} de um documento HTML.

	\vspace{0.3cm}
	Estrutura e significado ficam no HTML. Aparência fica no CSS. Comportamento
	fica no JavaScript, na segunda parte do semestre.
\end{frame}

\begin{frame}[fragile]{Onde o CSS entra: arquivo externo}
	A forma usada na disciplina e em produção:

	\begin{minted}{html}
<head>
  <meta charset="utf-8">
  <title>Catálogo</title>
  <link rel="stylesheet" href="css/estilo.css">
</head>
  \end{minted}

	\vspace{0.3cm}
	\begin{itemize}
		\item o arquivo \texttt{.css} não tem nenhuma tag HTML dentro;
		\item o \texttt{href} segue as regras de caminho relativo já conhecidas;
		\item um arquivo atende o site inteiro, e fica em cache.
	\end{itemize}
\end{frame}

\begin{frame}[fragile]{As outras duas formas}
	Elemento \texttt{<style>} no documento, para teste rápido:

	\begin{minted}{html}
<style>
  p { color: #333333; }
</style>
  \end{minted}

	\vspace{0.2cm}
	Atributo \texttt{style} na tag:

	\begin{minted}{html}
<p style="color: red;">Texto</p>
  \end{minted}

	\vspace{0.2cm}
	O atributo mistura apresentação e estrutura, obriga a repetir a declaração em
	cada elemento e vence qualquer seletor da folha. Voltaremos a ele na
	especificidade.
\end{frame}

\begin{frame}[fragile]{Anatomia de uma regra}
	\begin{minted}{css}
/* isto é um comentário, e é a única forma de comentar */
h1 {
  color: #1a3d6d;
  font-size: 2rem;
}
  \end{minted}

	\vspace{0.2cm}
	\begin{itemize}
		\item \texttt{h1}: o \textbf{seletor}, diz a quais elementos a regra se aplica;
		\item entre chaves: o \textbf{bloco de declarações};
		\item \texttt{color: \#1a3d6d;}: uma \textbf{declaração}, com
		      \textbf{propriedade}, dois pontos, \textbf{valor} e ponto e vírgula.
	\end{itemize}
\end{frame}

\begin{frame}{CSS não avisa quando você erra}
	\begin{alertblock}{Falha silenciosa}
		Propriedade inexistente, valor inválido, chave faltando: o navegador
		descarta o que não entendeu e segue adiante. Nenhuma mensagem aparece na
		tela.
	\end{alertblock}

	\vspace{0.5cm}
	Consequências para o dia a dia:

	\begin{itemize}
		\item ``não funcionou'' quase nunca significa que o navegador está errado;
		\item o painel \textbf{Styles} do inspetor mostra o que foi aplicado e o que
		      foi descartado;
		\item o validador do W3C aponta o nome digitado errado.
	\end{itemize}
\end{frame}

% ==========================================
\section{Seletores}
% ==========================================

\begin{frame}[fragile]{Tipo, classe e id}
	\begin{minted}{css}
/* tipo: todos os parágrafos */
p { line-height: 1.6; }

/* classe: todo elemento com class="destaque" */
.destaque { background-color: #fff3cd; }

/* id: o único elemento com id="rodape" */
#rodape { font-size: 0.875rem; }
  \end{minted}

	\vspace{0.2cm}
	A classe se repete à vontade no documento. O \texttt{id} é único.
\end{frame}

\begin{frame}[fragile]{Universal, agrupamento e relação}
	\begin{minted}{css}
/* universal: qualquer elemento */
* { box-sizing: border-box; }

/* agrupamento: a mesma declaração para vários seletores */
h1, h2, h3 { font-family: Georgia, serif; }

/* descendente: um <a> em qualquer nível dentro do <nav> */
nav a { text-decoration: none; }

/* filho direto: um <li> filho imediato do <ul> */
ul > li { margin-bottom: 0.5rem; }
  \end{minted}

	\vspace{0.2cm}
	O espaço entre dois seletores significa ``dentro de''.
\end{frame}

\begin{frame}[fragile]{Atributo, pseudo-classe e pseudo-elemento}
	\begin{minted}{css}
/* atributo */
input[type="text"] { border: 1px solid #999999; }

/* pseudo-classe: estado do elemento */
a:hover { text-decoration: underline; }
a:focus-visible { outline: 3px solid #1a73e8; }

/* pseudo-elemento: parte do elemento, sem tag própria */
p::first-line { font-variant: small-caps; }
  \end{minted}

	\vspace{0.2cm}
	Pseudo-classe: dois pontos, descreve estado. Pseudo-elemento: dois sinais de
	dois pontos, alcança um pedaço do elemento.
\end{frame}

\begin{frame}[fragile]{O contorno de foco não se apaga}
	\begin{alertblock}{Acessibilidade}
		Quem navega por teclado enxerga onde está pelo contorno de foco.
		\texttt{outline: none} deixa a pessoa perdida na página.
	\end{alertblock}

	\vspace{0.4cm}
	Se o contorno padrão não combina com o visual, troque a aparência dele:

	\begin{minted}{css}
a:focus-visible {
  outline: 3px solid #1a73e8;
  outline-offset: 2px;
}
  \end{minted}
\end{frame}

\begin{frame}{Classe ou id?}
	\begin{columns}[T]
		\column{0.48\textwidth}
		\textbf{Classe}
		\begin{itemize}
			\item reutilizável;
			\item quantos elementos você quiser;
			\item escolha padrão para estilo.
		\end{itemize}

		\column{0.48\textwidth}
		\textbf{\texttt{id}}
		\begin{itemize}
			\item único no documento;
			\item alvo de âncora (\texttt{href="\#rodape"});
			\item usado pelo JavaScript.
		\end{itemize}
	\end{columns}

	\vspace{0.6cm}
	Para estilo, prefira classe. O \texttt{id} tem especificidade alta demais, e a
	regra presa a ele fica difícil de sobrescrever depois. O próximo bloco mostra
	isso em números.
\end{frame}

% ==========================================
\section{Cascata}
% ==========================================

\begin{frame}[fragile]{Herança}
	\begin{minted}{css}
body {
  font-family: Verdana, sans-serif;
  color: #222222;
}
  \end{minted}

	\vspace{0.3cm}
	Todo o texto da página aparece em Verdana e cinza escuro, sem nenhuma regra
	para \texttt{p}, \texttt{li} ou \texttt{td}.

	\vspace{0.3cm}
	Algumas propriedades passam do elemento para os descendentes. Outras não.
\end{frame}

\begin{frame}{O que herda e o que não herda}
	\begin{center}
		\begin{tabular}{ll}
			\toprule
			\textbf{Herdam}      & \textbf{Não herdam}             \\
			\midrule
			\texttt{color}       & \texttt{margin}                 \\
			\texttt{font-family} & \texttt{padding}                \\
			\texttt{font-size}   & \texttt{border}                 \\
			\texttt{line-height} & \texttt{background}             \\
			\texttt{text-align}  & \texttt{width}, \texttt{height} \\
			\bottomrule
		\end{tabular}
	\end{center}

	\vspace{0.4cm}
	Propriedades de texto são herdadas. Propriedades de caixa não são, e é o
	comportamento desejado: um \texttt{padding} no \texttt{body} não vira
	\texttt{padding} em cada parágrafo.

	\vspace{0.2cm}
	A página da propriedade na MDN traz a linha \emph{Herdada: sim} ou
	\emph{não}.
\end{frame}

\begin{frame}{A cascata resolve conflitos em três passos}
	Duas regras dizem coisas diferentes sobre a mesma propriedade do mesmo
	elemento. O navegador decide nesta ordem:

	\begin{enumerate}
		\item \textbf{Origem e importância}: o autor da página vence o navegador;
		      \texttt{!important} sobe acima de todas as declarações normais;
		\item \textbf{Especificidade}: vence o seletor mais específico;
		\item \textbf{Ordem de aparição}: no empate, vence a última.
	\end{enumerate}

	\vspace{0.4cm}
	O passo 3 explica por que a ordem das regras no arquivo importa, e por que a
	folha de um framework vem antes da sua no \texttt{<head>}.
\end{frame}

\begin{frame}{Especificidade: contagem em três casas}
	\begin{center}
		\begin{tabular}{lccc}
			\toprule
			\textbf{Seletor}                  & \textbf{A} & \textbf{B} & \textbf{C} \\
			\midrule
			\texttt{p}                        & 0          & 0          & 1          \\
			\texttt{.texto}                   & 0          & 1          & 0          \\
			\texttt{p.texto}                  & 0          & 1          & 1          \\
			\texttt{.artigo p.destaque}       & 0          & 2          & 1          \\
			\texttt{\#noticia p}              & 1          & 0          & 1          \\
			\texttt{input[type="text"]:focus} & 0          & 2          & 1          \\
			\bottomrule
		\end{tabular}
	\end{center}

	\vspace{0.3cm}
	\textbf{A}: \texttt{id}. \textbf{B}: classes, atributos e pseudo-classes.
	\textbf{C}: tipos e pseudo-elementos. O \texttt{*} não conta nada.

	\vspace{0.2cm}
	Compara-se da esquerda para a direita: um \texttt{id} vence qualquer
	quantidade de classes.
\end{frame}

\begin{frame}[fragile]{Qual cor vence?}
	\begin{minted}{html}
<div class="artigo" id="noticia">
  <p class="texto destaque">O sistema entra em produção.</p>
</div>
  \end{minted}

	\begin{minted}{css}
p                   { color: blue; }
.texto              { color: green; }
.destaque           { color: orange; }
.artigo p.destaque  { color: purple; }
#noticia p          { color: red; }
  \end{minted}
\end{frame}

\begin{frame}{Vermelho}
	\texttt{\#noticia p} tem A igual a 1, e nenhuma quantidade de classes alcança
	isso.

	\vspace{0.5cm}
	Duas observações sobre o mesmo exemplo:

	\begin{itemize}
		\item entre \texttt{.texto} e \texttt{.destaque}, que empatam em (0,1,0),
		      venceria \texttt{.destaque}, por vir depois no arquivo;
		\item a ordem das classes dentro do atributo \texttt{class} do HTML não
		      influencia nada.
	\end{itemize}
\end{frame}

\begin{frame}{Acima de tudo: o atributo style, e o !important}
	O atributo \texttt{style} na tag vence qualquer seletor. É por isso que o
	estilo \emph{inline} atrapalha: ele tira de você a chance de corrigir aquela
	declaração pela folha.

	\vspace{0.4cm}
	\begin{alertblock}{Sobre o \texttt{!important}}
		Ele funciona, e cria um problema maior do que resolveu: a única forma de
		vencer um \texttt{!important} depois é outro \texttt{!important}. Considere-o
		proibido nesta disciplina.
	\end{alertblock}

	\vspace{0.3cm}
	Quando der vontade de usá-lo, o que falta é descobrir por que a sua regra
	perdeu. Faremos isso no inspetor.
\end{frame}

\begin{frame}{Seletor simples é seletor que dá para corrigir}
	\begin{block}{Regra de bolso}
		Escreva o seletor mais simples que resolve o problema.
	\end{block}

	\vspace{0.5cm}
	Uma folha construída sobre classes de um nível é fácil de ajustar depois.

	\vspace{0.3cm}
	Uma folha cheia de \texttt{\#menu .lista li a span} obriga cada correção
	futura a ser mais específica que a anterior, e o arquivo cresce em espiral.
\end{frame}

% ==========================================
\section{Caixa}
% ==========================================

\begin{frame}[fragile]{Todo elemento é uma caixa}
	\begin{minted}{text}
+-------------------------------------------+
|                  margin                   |
|   +-----------------------------------+   |
|   |              border               |   |
|   |   +---------------------------+   |   |
|   |   |         padding           |   |   |
|   |   |   +-------------------+   |   |   |
|   |   |   |     content       |   |   |   |
|   |   |   +-------------------+   |   |   |
|   |   +---------------------------+   |   |
|   +-----------------------------------+   |
+-------------------------------------------+
  \end{minted}

	\vspace{0.2cm}
	\texttt{padding} recebe a cor de fundo. \texttt{margin} é sempre transparente.
\end{frame}

\begin{frame}[fragile]{Formas abreviadas, no sentido horário}
	\begin{minted}{css}
padding: 10px;                /* os quatro lados */
padding: 10px 20px;           /* topo e base | direita e esquerda */
padding: 10px 20px 5px;       /* topo | direita e esquerda | base */
padding: 10px 20px 5px 0;     /* topo | direita | base | esquerda */
padding-left: 20px;           /* um lado só */

border: 5px solid #000000;    /* espessura, estilo e cor */
  \end{minted}

	\vspace{0.2cm}
	As mesmas formas valem para \texttt{margin} e para \texttt{border-width}.
\end{frame}

\begin{frame}[fragile]{A largura ocupada pela caixa}
	\begin{minted}{css}
.caixa {
  width: 300px;
  padding: 30px;
  border: 5px solid #000000;
}
  \end{minted}

	\vspace{0.3cm}
	Largura ocupada na tela:

	\begin{center}
		\large 300 + 30 + 30 + 5 + 5 = \textbf{370px}
	\end{center}

	\vspace{0.3cm}
	\texttt{width} dimensiona apenas o \texttt{content}. Duas colunas de
	\texttt{width: 50\%} com \texttt{padding} somam mais que a largura disponível,
	e a segunda desce para baixo da primeira.
\end{frame}

\begin{frame}[fragile]{box-sizing: border-box}
	\begin{minted}{css}
.caixa {
  box-sizing: border-box;
  width: 300px;
  padding: 30px;
  border: 5px solid #000000;
}
  \end{minted}

	\vspace{0.3cm}
	Agora a caixa ocupa \textbf{300px}, e o conteúdo é espremido para
	$300 - 60 - 10 = 230$px.

	\vspace{0.3cm}
	Por isso a linha abaixo abre praticamente toda folha de estilo moderna:

	\begin{minted}{css}
* { box-sizing: border-box; }
  \end{minted}
\end{frame}

\begin{frame}{Colapso de margens verticais}
	Primeira caixa com \texttt{margin-bottom: 30px}. Segunda com
	\texttt{margin-top: 20px}. O espaço entre elas é de \textbf{30px}.

	\vspace{0.4cm}
	Margens verticais adjacentes se fundem, e a maior prevalece.

	\begin{itemize}
		\item acontece também entre um elemento e o pai, quando não há borda nem
		      padding separando os dois;
		\item margens horizontais nunca colapsam;
		\item dentro de Flexbox e Grid não há colapso, o que é um dos motivos de
		      esses modelos deixarem o espaçamento previsível.
	\end{itemize}

	\vspace{0.3cm}
	Saída prática: espaçar sempre para o mesmo lado, usando só a margem de baixo
	entre irmãos.
\end{frame}

\begin{frame}{display}
	\begin{center}
		\small
		\begin{tabular}{lccc}
			\toprule
			\textbf{Valor}        & \textbf{Ocupa a linha}                        & \textbf{Aceita \texttt{width}} &
			\textbf{Margem vertical}                                                                                             \\
			\midrule
			\texttt{block}        & sim                                           & sim                            & empurra     \\
			\texttt{inline}       & não                                           & não                            & não empurra \\
			\texttt{inline-block} & não                                           & sim                            & empurra     \\
			\texttt{none}         & \multicolumn{3}{c}{o elemento some da página}                                                \\
			\bottomrule
		\end{tabular}
	\end{center}

	\vspace{0.4cm}
	\texttt{<div>}, \texttt{<p>}, \texttt{<h1>} e \texttt{<section>} são
	\texttt{block}. \texttt{<a>}, \texttt{<span>}, \texttt{<strong>} e
	\texttt{<img>} são \texttt{inline}.

	\vspace{0.2cm}
	\texttt{display: none} também tira o elemento do leitor de tela.
\end{frame}

\begin{frame}[fragile]{Link vira botão com inline-block}
	\begin{minted}{css}
.botao {
  display: inline-block;
  padding: 0.75rem 1.5rem;
  background-color: #1a3d6d;
  color: #ffffff;
  text-decoration: none;
}
  \end{minted}

	\vspace{0.3cm}
	Sem o \texttt{inline-block}, o \texttt{padding} vertical não afasta as linhas
	vizinhas e \texttt{width} é ignorado, porque \texttt{<a>} é \texttt{inline}.
\end{frame}

% ==========================================
\section{Unidades e cores}
% ==========================================

\begin{frame}[fragile]{Unidades de medida}
	\begin{minted}{css}
.exemplo {
  width: 300px;    /* pixel CSS, medida absoluta */
  width: 50%;      /* relativo ao elemento pai */
  font-size: 1rem; /* relativo à fonte do <html> */
  padding: 0.5em;  /* relativo à fonte do próprio elemento */
  height: 100vh;   /* 100% da altura da janela */
  width: 80vw;     /* 80% da largura da janela */
}
  \end{minted}
\end{frame}

\begin{frame}{Quando usar cada uma}
	\begin{itemize}
		\item \texttt{px}: bordas e detalhes que não devem escalar;
		\item \texttt{rem}: texto e espaçamento. Apoia-se na fonte do \texttt{<html>},
		      que é a preferência configurada por quem lê;
		\item \texttt{em}: se apoia na fonte do próprio elemento, e por isso acumula
		      em elementos aninhados. Útil no padding de botão;
		\item \texttt{\%}: depende da propriedade. Em \texttt{width}, é a largura do
		      pai;
		\item \texttt{vw} e \texttt{vh}: centésimos da janela, para seções de tela
		      cheia.
	\end{itemize}
\end{frame}

\begin{frame}[fragile]{Não fixe o tamanho de fonte do html}
	\begin{alertblock}{Linha a evitar}
		\texttt{html \{ font-size: 16px; \}}
	\end{alertblock}

	\vspace{0.4cm}
	Anula o ajuste de tamanho de fonte de
	quem enxerga pouco, uma das configurações de acessibilidade mais usadas.

	\vspace{0.3cm}
	O padrão do \texttt{<html>} já é a preferência da pessoa. Para deixar isso
	explícito:

	\begin{minted}{css}
html { font-size: 100%; }
  \end{minted}
\end{frame}

\begin{frame}[fragile]{Notações de cor}
	\begin{minted}{css}
color: red;                     /* palavra-chave */
color: #1a3d6d;                 /* hexadecimal RR GG BB */
color: #fff;                    /* forma curta de #ffffff */
color: rgb(26, 61, 109);        /* 0 a 255 em cada canal */
color: rgba(26, 61, 109, 0.5);  /* o quarto valor é a opacidade */
color: hsl(212, 61%, 26%);      /* matiz, saturação, luminosidade */
  \end{minted}

	\vspace{0.2cm}
	O hexadecimal é o que aparece em qualquer paleta pronta. O \texttt{hsl()} é o
	mais fácil de ajustar à mão: para o azul mais escuro do \texttt{:hover}, basta
	reduzir o terceiro valor.
\end{frame}

\begin{frame}{Contraste mínimo entre texto e fundo}
	Texto cinza claro sobre fundo branco fica ilegível para quem tem baixa visão,
	para quem usa a tela sob luz forte e para parte das pessoas com daltonismo.

	\vspace{0.4cm}
	A WCAG define a razão mínima de contraste entre texto e fundo:

	\begin{itemize}
		\item \textbf{4.5:1} para texto normal;
		\item \textbf{3:1} para texto grande, a partir de 24px, ou 18.66px em
		      negrito.
	\end{itemize}

	\vspace{0.3cm}
	Confira em \url{https://webaim.org/resources/contrastchecker/}, colando os dois
	hexadecimais. O inspetor mostra a mesma razão ao lado da cor.

	\vspace{0.2cm}
	\small ODS 4 e ODS 10: contraste insuficiente exclui parte do público sem que
	o autor perceba: no monitor dele, o texto está legível.
\end{frame}

% ==========================================
\section{Tipografia}
% ==========================================

\begin{frame}[fragile]{font-family recebe uma lista}
	\begin{minted}{css}
body {
  font-family: "Segoe UI", Roboto, Helvetica, Arial, sans-serif;
}
  \end{minted}

	\vspace{0.3cm}
	O navegador tenta a primeira, e passa à seguinte se ela não estiver instalada
	na máquina.

	\begin{itemize}
		\item a lista termina em família genérica: \texttt{sans-serif},
		      \texttt{serif} ou \texttt{monospace};
		\item nomes com espaço vão entre aspas.
	\end{itemize}
\end{frame}

\begin{frame}[fragile]{Tamanho, peso e altura de linha}
	\begin{minted}{css}
body { font-size: 1rem; line-height: 1.6; }
h1   { font-size: 2.5rem; font-weight: 700; }
  \end{minted}

	\vspace{0.3cm}
	\texttt{line-height} \textbf{sem unidade}: o valor é herdado como fator, e cada
	elemento o multiplica pelo próprio tamanho de fonte. Entre 1.4 e 1.6 para texto
	corrido.

	\vspace{0.3cm}
	\texttt{font-weight}: \texttt{normal} é 400, \texttt{bold} é 700, e os
	múltiplos de 100 valem quando a fonte tem esses cortes.
\end{frame}

\begin{frame}[fragile]{Fonte externa do Google Fonts}
	No \texttt{<head>}, antes da sua folha:

	\begin{minted}{html}
<link rel="preconnect" href="https://fonts.googleapis.com">
<link href="https://fonts.googleapis.com/css2?family=Roboto:wght@400;700&display=swap" rel="stylesheet">
  \end{minted}

	No CSS, mantendo a pilha de reserva:

	\begin{minted}{css}
body { font-family: Roboto, Arial, sans-serif; }
  \end{minted}

	\vspace{0.2cm}
	Cada peso pedido é um arquivo a mais para baixar. Peça dois, no máximo três.
\end{frame}

% ==========================================
\section{Prática}
% ==========================================

\begin{frame}{Roteiro da prática}
	Trabalhando sobre as páginas da tarefa de HTML:

	\begin{enumerate}
		\item criar \texttt{css/estilo.css} e ligá-lo pelo \texttt{<link>};
		\item descobrir no inspetor por que uma regra não pegou;
		\item medir a caixa no diagrama do painel \textbf{Computed};
		\item verificar contraste e validar a folha no W3C.
	\end{enumerate}

	\vspace{0.4cm}
	O texto completo do roteiro está no material da aula, na parte prática.
\end{frame}

\begin{frame}[fragile]{Por que a regra não pegou}
	Com \texttt{class="menu"} no \texttt{<nav>}, e estas três regras nesta ordem:

	\begin{minted}{css}
nav a   { color: #1a3d6d; }
a       { color: #cc0000; }
.menu a { color: #008000; }
  \end{minted}

	\vspace{0.2cm}
	Decida a cor antes de recarregar. Depois, no painel \textbf{Styles}:

	\begin{itemize}
		\item as declarações perdedoras aparecem \textbf{riscadas};
		\item ao lado de cada regra, o arquivo e a linha onde ela está;
		\item a aba \textbf{Computed} mostra o valor final já resolvido.
	\end{itemize}
\end{frame}

\begin{frame}[fragile]{Medir a caixa}
	\begin{minted}{css}
.caixa {
  width: 300px; padding: 30px;
  border: 5px solid #000000;
  background-color: #ffd400;
}
.corrigida { box-sizing: border-box; }
  \end{minted}

	\vspace{0.2cm}
	No fim do painel \textbf{Computed}, o diagrama do modelo de caixa traz os
	quatro valores. Confirme: \textbf{370px} na primeira, \textbf{300px} na
	segunda.
\end{frame}

\begin{frame}{Exercícios}
	\begin{block}{Entrega}
		Valem nota, no item \emph{Exercícios em sala}. O \textbf{prazo} está na
		atividade correspondente na \textbf{UFPR Virtual}, e a entrega é por
		\texttt{push} no GitLab.
	\end{block}

	\vspace{0.4cm}
	\begin{itemize}
		\item enunciado no \texttt{README.md} do repositório-modelo, com o endereço
		      na UFPR Virtual;
		\item individual ou em dupla, com \emph{fork} dentro do grupo;
		\item atividade avaliativa: sem IA generativa para produzir o código;
		\item a folha que sair daqui recebe o layout no próximo encontro e vai na
		      Entrega 1 do trabalho.
	\end{itemize}
\end{frame}

\begin{frame}{O que ainda falta para montar um layout}
	Você controla a aparência de cada caixa. Falta posicionar as caixas umas em
	relação às outras: menu na horizontal, três cartões lado a lado, coluna lateral
	acompanhando a altura do conteúdo.

	\vspace{0.4cm}
	Próximo encontro: Flexbox, Grid e \emph{media queries}.

	\vspace{0.3cm}
	Os fundamentos de hoje continuam valendo lá dentro: a caixa que o Flexbox
	posiciona é a mesma caixa desta aula.
\end{frame}

\begin{frame}{Resumo}
	\begin{itemize}
		\item toda página já vem estilizada pela folha do navegador;
		\item conflito se resolve por origem, especificidade e ordem, nessa sequência;
		\item um \texttt{id} vence qualquer quantidade de classes; prefira classe;
		\item propriedades de texto são herdadas, as de caixa não;
		\item \texttt{width} dimensiona só o conteúdo, até você escrever
		      \texttt{box-sizing: border-box};
		\item margens verticais adjacentes colapsam;
		\item \texttt{rem} para texto, e sem fixar a fonte do \texttt{<html>};
		\item contraste de 4.5:1 é requisito;
		\item CSS inválido não gera erro na tela: a resposta está no inspetor.
	\end{itemize}
\end{frame}

\end{document}
